Mit 1~Zoll $= \zoO$ ist der Durchmesser $d = \dZO\cdot\zo = \dplP{-2}{3}$, die Frequenz $f = \fUO = \fUsP{0}{2}$.
\begin{abcliste}
\abc Ein Punkt am Rand legt pro Umdrehung den Umfang $\pi d$ zurück:
\begin{align}
 v &= \vplF \\
   &= \pi\cdot\dpl\cdot\fUs \\
   &= \vpl \approx \result{\vplP{0}{2}}
\end{align}

\abc Im ungünstigsten Fall ist der gesuchte Punkt gerade am Lesekopf vorbeigelaufen und muss eine ganze Umdrehung abwarten:
\begin{align}
 T &= \TplF = \frac{1}{\fUs} = \Tpl \approx \result{\TplP{-3}{2}}
\end{align}
Alle Punkte der Platte haben dieselbe Umlaufzeit (gleiche Winkelgeschwindigkeit). Die maximale Zugriffszeit hängt deshalb \result{\text{nicht}} von der Entfernung zur Plattenmitte ab; im Mittel wartet man eine halbe Umdrehung.
\end{abcliste}
